\documentclass[11pt]{article}
\usepackage[T1]{fontenc}
\usepackage[utf8]{inputenc}
\usepackage[english]{babel}
\usepackage[margin=1in]{geometry}
\usepackage{amsmath,amssymb}
\usepackage{enumitem}
\usepackage{microtype}

\ifdefined\pdfcatalog
  \pdfcatalog{/Lang (en-US)}
\fi

\setlength{\parindent}{0pt}
\setlength{\parskip}{0.55em}
\setlength{\emergencystretch}{2em}
\setlist{nosep,leftmargin=*}

\newcommand{\findinglabel}[1]{\textbf{#1}}
\newcommand{\classification}[1]{\textit{Classification: #1.}}

\title{Technical Review of ``Early-Stopping Activation Steering for Factual Preference Alignment in Vietnamese Medical Language Models''}
\author{Manuscript review}
\date{August 24, 2026}

\begin{document}
\selectlanguage{english}
\maketitle

\section*{Overall assessment}

The revised manuscript adds two potentially valuable components: a four-condition natural-completion comparison and Layer~8 activation-norm tracking.  It also corrects the absolute-dose formulas, random-vector indexing, category-table column count, and some earlier wording.  The bounded 200-token McNemar result remains internally consistent and is still the clearest supported result.

The new material nevertheless introduces major contradictions.  The natural-completion prose reports an obsolete 64.80\%--65.80\% comparison, while its table and the Abstract report different baseline and steering values.  The table also shows that every steering condition has lower BERTScore and higher repetition than the baseline, despite higher reference preference; this does not support unqualified ``outperforming'' or repetition-prevention language.  The activation-norm diagnostic establishes, at most, a difference in scalar norm after different interventions.  It does not establish representation saturation, departure from a hidden-state manifold, or a causal link to output quality.  Placebo, RAG, latency, dataset, clinical-validity, and reproducibility limitations also remain.  The manuscript therefore still requires major revision.

\section*{Major technical findings}

\subsection*{1. The primary natural-completion results are internally contradictory}

\classification{Definite data-reporting error}

\findinglabel{Location.} Abstract; Primary Natural Completion Evaluation; Table~\texttt{tab:long\_gen}; Discussion; Conclusion.

\findinglabel{Problem.} The paragraph preceding Table~\texttt{tab:long\_gen} reports a baseline of 64.80\% (324/500), early-stop steering of 65.80\% (329/500), BERTScore $0.6804\rightarrow0.6834$, and $p=0.4583$.  The table instead reports 65.40\% for baseline, 68.60\% for continuous steering, 67.80\% for linear decay, and 70.60\% for hard cutoff, with BERTScores 0.6779, 0.6705, 0.6732, and 0.6695.  These percentages imply counts of 327, 343, 339, and 353 out of 500, respectively, not the counts in the prose.  The Abstract follows the new table, whereas the Discussion still reports the obsolete $+1.00$-point trend.  The text repeatedly states a 100.00\% EOS hit rate for the natural-completion experiment, but the table gives 99.80\% for continuous steering and linear decay.

\findinglabel{Why it matters.} The paper's primary evaluation cannot be interpreted or reproduced until one set of outputs is identified as authoritative.  The discrepancy changes the effect size, the best schedule, the statistical conclusion, and the BERTScore direction.

\findinglabel{Suggested fix.} Regenerate the Abstract, Results, Discussion, and Conclusion from one version-controlled analysis output.  Report raw counts for all four conditions, the full paired transition tables, paired effect estimates and confidence intervals, exact test definitions, and finish-reason counts.  Remove every obsolete number before making any scientific interpretation.

\subsection*{2. The natural-completion table shows a metric trade-off, not general outperformance}

\classification{Definite interpretation error}

\findinglabel{Location.} Abstract; Table~\texttt{tab:long\_gen}; Empirical Representation Saturation Diagnostics; Discussion; Conclusion.

\findinglabel{Problem.} Relative to the table's baseline, hard cutoff increases RefPref by 5.20 points and linear decay by 2.40 points.  However, all steering schedules lower BERTScore and increase Rep-4:
\[
\begin{array}{c|ccc}
& \text{Baseline} & \text{Linear decay} & \text{Hard cutoff} \\
\hline
\text{BERTScore} & 0.6779 & 0.6732 & 0.6695 \\
\text{Rep-4 (\%)} & 4.10 & 5.37 & 5.35
\end{array}
\]
Continuous steering also has lower BERTScore (0.6705) and higher Rep-4 (4.38\%) than baseline.  Linear decay has the highest repetition rate of the four conditions.  Hard cutoff, not linear decay, has the highest reference-preference rate, while baseline has the best BERTScore and repetition.  No uncertainty is reported for any of these differences.

\findinglabel{Why it matters.} The new primary experiment does not support the claims that steering generally outperforms baseline, that linear decay is the best schedule, or that early stopping prevents repetitive output.  It shows a multi-objective trade-off whose clinical meaning is unknown.

\findinglabel{Suggested fix.} Report paired confidence intervals and prespecified comparisons for every primary metric.  Describe the result as increased reference preference accompanied by lower reference BERTScore and higher repetition unless corrected data show otherwise.  Select a primary schedule and outcome using validation data, state the selection rule, and evaluate that frozen choice once on the test set.  Do not use ``outperforming'' without naming the metric.

\subsection*{3. Activation-norm differences do not demonstrate representation saturation}

\classification{Unsupported mechanism claim and implementation ambiguity}

\findinglabel{Location.} Abstract; contribution~3; Early-Stopping Steering Hook \& Saturation Diagnostics; Empirical Representation Saturation Diagnostics.

\findinglabel{Problem.} The diagnostic reports selected values of $\|h_8^{(t)}\|_2$ and interprets a larger scalar norm under continuous steering as ``saturation'' and movement ``outside the unsteered activation manifold.''  Neither conclusion follows from the norm alone.  Saturation requires an operational definition, such as a nonlinear loss of responsiveness, clipping, directional collapse, or a measured degradation associated with the state change.  A scalar radius cannot identify whether a point is outside a high-dimensional manifold.  Moreover, once early steering is set to zero after $K=16$, a smaller direct perturbation at $t=50$ is expected by construction; this observation alone does not show stabilization or explain output behavior.

The manuscript does not define whether it averages per-token norms or takes the norm of an averaged activation, whether $h_l^{(t)}$ is measured before intervention or $\hat h_l^{(t)}$ after intervention, which token positions are modified, or how prompts that terminate before $t=50$ or $t=100$ are handled.  No variance, confidence interval, per-step sample size, trajectory figure, or paired relation to repetition/correctness is reported.  The hard-cutoff statement gives a drop from 56.28 to 52.75 without identifying the two decoding steps.

\findinglabel{Why it matters.} The activation diagnostic is used as direct evidence for the paper's central mechanism, but its current design and interpretation only establish unquantified differences in one summary statistic.

\findinglabel{Suggested fix.} Define saturation before analysis and report complete trajectories with per-step sample counts and uncertainty.  Distinguish pre-hook $h_l^{(t)}$ from post-hook $\hat h_l^{(t)}$, state the aggregation rule, and handle variable output lengths explicitly.  Measure cosine drift, projection along $v_{\text{steer}}$, covariance or distributional distance, response to additional perturbation, and associations with repetition and factual errors.  Replace ``outside the manifold'' and ``confirms saturation'' with ``elevated Layer~8 norm'' unless stronger diagnostics support those conclusions.

\subsection*{4. The central outcome remains a semantic proxy rather than clinical factual correctness}

\classification{Unsupported interpretation of a correctly disclosed proxy}

\findinglabel{Location.} Abstract; Introduction; contributions; Results; Conclusion; keywords.

\findinglabel{Problem.} The Methods correctly states that RefPref is a BERTScore proxy, not clinician-adjudicated correctness.  Elsewhere the manuscript calls it ``factual alignment,'' ``truthfulness,'' ``superior accuracy,'' and hallucination mitigation.  A generated answer can be closer to $y^+$ than to one synthetic $y^-$ while still containing a wrong dose, unit, negation, contraindication, or interaction mechanism.  The natural-completion table further shows that RefPref can increase while BERTScore to the reference decreases, emphasizing that the proxy requires careful interpretation.

\findinglabel{Why it matters.} Clinical safety cannot be inferred from a relative embedding comparison against two templates.  The strongest title-level and deployment claims therefore exceed the measured endpoint.

\findinglabel{Suggested fix.} Make blinded clinician-adjudicated factual correctness and harmful-error rate the primary endpoints, stratified by pregnancy, interaction, and dosage cases.  Validate RefPref against those labels and report sensitivity, specificity, agreement, and failure cases.  Until then, use ``reference-preference rate'' consistently and avoid presenting it as clinical accuracy or hallucination mitigation.

\subsection*{5. The random-direction controls support specificity but remain statistically overstated}

\classification{Unsupported strength of evidence}

\findinglabel{Location.} Abstract; contribution~4; directional-control Methods; Placebo Control Distribution; Table~\texttt{tab:baselines}; Conclusion.

\findinglabel{Problem.} The standardized distance
\[
\frac{77.20-56.50}{5.11}\approx4.05
\]
is not a calibrated hypothesis test or a $4\sigma$ confidence statement.  With 20 sampled directions and none reported above 63.00\%, the simplest corrected empirical randomization bound is approximately $1/(20+1)=0.0476$, assuming a prespecified exchangeable sampling procedure.  The individual direction results, seeds, BERTScore dispersion, and paired question-level uncertainty are absent.  The table gives one BERTScore value (0.6945) for the entire distribution without defining its aggregation.  Negative steering changes RefPref by $-10.20$ points whereas positive steering changes it by $+4.20$ points, so the result is not symmetric despite the stated purpose of the control.

\findinglabel{Why it matters.} Random and sign-flipped controls are useful evidence that direction matters, but they do not prove that the direction represents factuality or establish the stated Gaussian-tail interpretation.

\findinglabel{Suggested fix.} Release all seeds and per-direction, per-question outputs.  Report mean, sample SD, range, and confidence intervals for every metric, and use a prespecified randomization analysis or a hierarchical bootstrap over directions and questions.  Analyze the positive/negative asymmetry.  Describe the result as direction-specific evidence rather than proof or a calibrated $4\sigma$ effect.

\subsection*{6. The BM25 baseline is insufficiently specified and too weak for a general RAG comparison}

\classification{Likely weak baseline and unsupported causal attribution}

\findinglabel{Location.} Experimental Setup; Placebo Control Distribution \& Real BM25 RAG Results; Abstract; Conclusion.

\findinglabel{Problem.} Top-1 Recall@1 is only 13.00\%, but the manuscript does not define corpus construction, passage boundaries, Vietnamese tokenization, query text, BM25 implementation or parameters, relevance mapping, prompt template, or retrieval-failure handling.  It attributes the 68.60\% RefPref result to retrieval noise without an oracle-context control or retrieval-conditioned error analysis.  No paired interval or test compares BM25 RAG with baseline or steering.  A 13\% recall may reflect an underdeveloped retriever rather than a general limitation of RAG.

\findinglabel{Why it matters.} The claim that steering is more accurate and efficient than RAG is a headline result, but it currently compares against one unreproducible and apparently weak retrieval configuration.

\findinglabel{Suggested fix.} Fully specify and release the retrieval pipeline.  Report Recall@$k$ for several values of $k$, stronger lexical or dense retrievers, a gold-context/oracle condition, and generation results conditional on retrieval success.  Use paired tests on identical questions.  Describe the current result as one BM25 configuration rather than a general RAG comparison.

\subsection*{7. Latency and memory claims remain unsupported}

\classification{Definite systems-reporting insufficiency}

\findinglabel{Location.} Experimental Setup; Tables~\texttt{tab:baselines} and \texttt{tab:equal\_alpha\_ablation}; Abstract; Conclusion.

\findinglabel{Problem.} Every non-RAG condition, including all 12 ablation rows, negative steering, and the 20-direction aggregate, is reported as exactly 7.32~s.  No timing repetitions, variance, warm-up, synchronization, batch size, GPU count, timing boundary, output-length control, or peak-memory measurement is supplied.  ``No measurable generation latency overhead'' is therefore not demonstrated by an equivalence test or uncertainty interval.  ``Zero-context-memory advantage'' is also imprecise: steering avoids retrieved context and index memory, but still uses a steering vector, hooks, activations, and the model's ordinary KV cache.  It is unclear whether the 13.70~s RAG measurement includes retrieval, prefill, decoding, or all three.

\findinglabel{Why it matters.} Systems claims cannot be inferred from one repeated rounded number, and the 87.1\% comparison is not interpretable without common timing boundaries.

\findinglabel{Suggested fix.} Run synchronized repeated timings after warm-up under identical hardware, batch size, prompts, and output-length controls.  Report mean, dispersion, trial count, prefill/decode/retrieval components, throughput, and peak CPU/GPU memory.  Use ``no measurable overhead'' only if an equivalence interval supports a prespecified negligible margin.

\subsection*{8. The ablation lacks uncertainty and the matched-dose sign definition is inconsistent}

\classification{Unsupported comparative claim and definite formula inconsistency}

\findinglabel{Location.} cumulative-dose equations; Controlled Factorial Ablation Matrix; Discussion.

\findinglabel{Problem.} The manuscript acknowledges that equal-$\alpha_0$ schedules have unequal $D_1$ and that the matched-$D_1$ control does not match temporal placement or $D_2=\sum_t\alpha(t)^2$.  Nevertheless, raw winners are emphasized for BERTScore differences as small as 0.0005--0.0018, without paired confidence intervals or multiplicity handling.  In addition, $D_{\text{decay}}$ is defined as a nonnegative absolute dose, but
\[
\alpha_{\text{match}}=D_{\text{decay}}/T
=\frac{K+1}{2T}\alpha_0
\]
cannot hold for negative $\alpha_0$: the first expression is positive and the second is negative.  The claim $\sum_t\alpha_{\text{match}}=D_{\text{decay}}$ therefore mixes signed and absolute quantities.

\findinglabel{Why it matters.} The sign mismatch makes the implementation-facing definition mathematically inconsistent outside the positive ablation rows, and the small untested rankings do not establish schedule superiority.

\findinglabel{Suggested fix.} Define a signed matched coefficient as
\[
\alpha_{\text{match}}
=\operatorname{sgn}(\alpha_0)\frac{D_1}{T},
\qquad
\sum_t|\alpha_{\text{match}}|=D_1,
\]
or restrict the definition explicitly to $\alpha_0>0$.  Report paired schedule contrasts with 95\% CIs and an exploratory multiplicity policy.  Avoid selecting a universal best schedule when different metrics prefer different rows.

\subsection*{9. Vector extraction and layer-selection evidence do not establish a transferable factual direction}

\classification{Likely representation confound and unsupported inference}

\findinglabel{Location.} Contrastive Vector Estimation; Early-Stopping Steering Hook; Layer-Wise Subspace Probing.

\findinglabel{Problem.} The direction is extracted from the final token of the completed sequence $x_i\oplus y_i^{\pm}$ but injected into the first 16 generated tokens.  Final states may encode answer length, punctuation, response templates, or synthetic-error artifacts rather than a context-invariant factual property.  The layer experiment is a steering-layer sweep, not a reported subspace probe.  Only the range 0.7040--0.7140 is supplied, without per-layer values or uncertainty, yet the paper infers that truthfulness is distributed across layers and that framework robustness is confirmed.

\findinglabel{Why it matters.} Proxy improvement does not establish the claimed representation semantics, and the extraction/injection temporal mismatch is an untested generalization assumption.

\findinglabel{Suggested fix.} Compare prompt-only, answer-initial, final-token, and pooled-response vectors; test independently authored errors and source/template-disjoint data; and report every layer's validation result with uncertainty.  Call the experiment a steering-layer sweep unless an actual trained probe or subspace analysis is performed.

\subsection*{10. Dataset, metric, statistical, and hook specifications remain incomplete}

\classification{Reproducibility limitation and likely leakage risk}

\findinglabel{Location.} benchmark construction; BERTScore criterion; Experimental Setup; statistical analysis; hook description.

\findinglabel{Problem.} ``Expert-structured'' remains undefined, with no annotator qualifications, review fraction, agreement, adjudication, counter-answer generation rules, quality checks, deduplication, source identifiers, or licensing.  Question-disjointness does not ensure drug/source-entry or perturbation-template disjointness.  The 500-question selection rule and seed are absent.  The paired interval method, Wilcoxon statistic, zero-difference handling, bootstrap details, tie counts, and multiplicity policy are not reported.  Rep-4 lacks a formula and tokenizer; Vietnamese ROUGE tokenization and full BERTScore settings are missing.  The exact Qwen module, layer-index convention, hook tensor positions, prompt-forward behavior, KV-cache handling, dtype/device, software versions, seeds, and GPU count are unspecified.

\findinglabel{Why it matters.} These omissions can materially change the vector, metrics, test statistics, and timing, and source/template leakage could inflate every reported result.

\findinglabel{Suggested fix.} Add a dataset card, source/template-grouped split manifests, expert-validation protocol, evaluation-subset seed, complete metric formulas/configurations, statistical software calls, environment details, and executable hook pseudocode.  Release per-question outputs and analysis code.

\section*{Equation, notation, sign, branch, and dimensional audit}

The direction $\mu_l^+-\mu_l^-$ and addition of $+\alpha(t)v_{\text{steer}}^{(l)}$ are sign-consistent with the intended positive intervention.  Unit normalization makes the vector dimensionally compatible with the hidden state when $\alpha(t)$ is an activation-scale coefficient.  The linear schedule and positive-$\alpha_0$ values of $D_{\text{continuous}}$, $D_{\text{cutoff}}$, and $D_{\text{decay}}=(K+1)|\alpha_0|/2$ are arithmetically correct.  No inverse-trigonometric functions, coordinate transforms, or branch/quadrant choices occur.

The concrete issues are the signed/absolute mismatch in $\alpha_{\text{match}}$, ambiguity between pre-intervention $h_l^{(t)}$ and post-intervention $\hat h_l^{(t)}$ in the norm diagnostic, and overload of $N$ for training samples, evaluation prompts, category sizes, and random directions.  The discrete token range should be written $t\in\{1,\ldots,100\}$ rather than $t\in[1,100]$.  Terminology also drifts between ``truthfulness,'' ``factual preference,'' and the observed RefPref proxy, while ``saturation,'' ``magnitude inflation,'' and ``outside the manifold'' are used as if equivalent.

\section*{Meaningful editorial, compilation, and reference findings}

\subsection*{The source compiles but contains visible Markdown and severe table overflow}

\findinglabel{Location.} Experimental Setup; Tables~\texttt{tab:long\_gen} and \texttt{tab:mcnemar}; clean pdfLaTeX output.

\findinglabel{Problem.} The Experimental Setup uses Markdown syntax such as \verb|**Dual-Evaluation Paradigm**| and plain numbered lines rather than LaTeX environments; the compiled PDF visibly prints the asterisks.  After two pdfLaTeX passes, the source builds to six pages, but the log reports overfull boxes of approximately 61.44~pt for the natural-completion table and 5.54~pt for the McNemar table.

\findinglabel{Why it matters.} The visible markup is unprofessional, and the natural-completion table can intrude substantially into the adjacent IEEE column.

\findinglabel{Suggested fix.} Replace Markdown with \verb|\textbf{...}| and an \texttt{enumerate} environment.  Redesign or resize the natural-completion and McNemar tables while preserving legibility, then rebuild twice and inspect the final PDF and log.

\subsection*{Several citations do not directly support the attached claims}

\findinglabel{Location.} Introduction; bibliography entries \texttt{ref10}, \texttt{ref11}, and \texttt{ref12}.

\findinglabel{Problem.} \texttt{ref10} is the LoRA paper rather than direct evidence for broad SFT hallucination mitigation.  \texttt{ref11} concerns robustness to irrelevant retrieved context but does not establish that internal activation pathways cause evidence ignoring.  \texttt{ref12} addresses catastrophic forgetting generally rather than the claimed medical-language-model failure mode.

\findinglabel{Why it matters.} The manuscript blurs direct prior evidence, mechanistic hypotheses, and general background.

\findinglabel{Suggested fix.} Add direct sources for hallucination-focused SFT, language-model forgetting after domain adaptation, and evidence-use failures in RAG.  Present the internal-activation explanation as a hypothesis unless it is measured or directly cited.

\section*{Proofing sweep}

The bundled mechanical scan was run on both the source and the clean compiled PDF and returned no high-confidence regex hits.  Manual source, dense-equation, table, and bibliography checks identified the following bounded corrections:

\begin{itemize}
    \item \textbf{Abstract, contribution~4, and RAG Results:} change ``an non-oracle'' to ``a non-oracle.''
    \item \textbf{Contribution~2:} replace ``un-truncated'' with ``untruncated.''
    \item \textbf{Experimental Setup:} replace Markdown bold and plain numbered text with valid LaTeX formatting.
    \item \textbf{Saturation diagnostics:} write ``activation $\ell_2$ norm'' rather than the redundant phrase ``elevated magnitude norm,'' and state the missing time indices for the hard-cutoff change from 56.28 to 52.75.
    \item \textbf{Natural-completion table:} define the boldface convention explicitly because different rows are best for RefPref, BERTScore, EOS, and Rep-4.
\end{itemize}

The bibliography was checked for duplicate punctuation, malformed quotation punctuation, and obvious citation-format glitches; no additional mechanical defect was found.

\section*{Items requiring external verification}

\begin{itemize}
    \item Verify all formulary-derived references and synthetic counter-answers against the official source, current clinical guidance, qualified clinician review, and applicable reuse permissions.
    \item Validate the multilingual BERTScore configuration for Vietnamese clinical negation, quantities, units, contraindications, and interaction mechanisms before interpreting RefPref as factual alignment.
    \item Verify novelty against current work on token-dependent, scheduled, and dose-controlled activation steering; the present literature review does not establish novelty of linear-decay early stopping.
    \item Verify the characterization of Qwen2.5-7B-Instruct as a ``Small Language Model'' and the privacy/local-hospital deployment claims using explicit definitions, a threat model, and measured resource requirements.
    \item Verify the BM25 baseline against established Vietnamese retrieval configurations and stronger retrievers before making a general comparison with RAG.
\end{itemize}

\section*{Highest-priority revision sequence}

\begin{enumerate}
    \item Resolve the contradictory natural-completion results and regenerate every summary from one analysis output.
    \item Reframe the primary result as a metric trade-off and add paired uncertainty for all natural-completion conditions.
    \item Redesign the activation diagnostic so that it operationally tests saturation, includes uncertainty, and connects hidden states to output behavior.
    \item Add clinician-adjudicated factual and harmful-error endpoints, or consistently narrow claims to reference preference.
    \item Fully report the random-direction distribution, strengthen and document RAG, and rerun systems timing with uncertainty.
    \item Add schedule-ablation uncertainty, correct the matched-dose sign definition, and test vector extraction/layer alternatives.
    \item Complete dataset provenance, leakage controls, metric/hook specifications, LaTeX formatting, table layout, and citation fit.
\end{enumerate}

\end{document}